\documentclass[11pt,onecolumn]{IEEEtran}
\usepackage{amsmath,amssymb,amsfonts}
\usepackage{fontspec}
\usepackage{unicode-math}
\setmainfont{Libertinus Serif}
\setmathfont{Libertinus Math}
\usepackage{algorithmic}
\usepackage{graphicx}
\usepackage{textcomp}
\usepackage{xcolor}
\usepackage{booktabs}
\usepackage{float}

\begin{document}

\title{Master's Project}

\author{Ali Ashoori \\
Institute for Advanced Studies in Basic Sciences (IASBS), Zanjan, Iran}

\maketitle

\begin{abstract}
  Phase unwrapping is a mathematically ill-posed inverse problem critical to optical interferometry, synthetic aperture radar (SAR), and quantitative phase imaging. Classical methods, including branch-cut and minimum-norm algorithms, are notoriously fragile in the presence of intense optical noise and dense fringe discontinuities. In this proposal, we introduce a fully convolutional deep learning architecture designed for robust, direct absolute phase reconstruction from raw interferograms. We detail a modified U-Net architecture that explicitly integrates normalized spatial coordinate embeddings to resolve global scalar ambiguities. Furthermore, we propose a mathematically rigorous curriculum noise synthesis pipeline that progressively exposes the network to physically accurate optical aberrations, laser speckle, and quantum shot noise. By employing a multi-term loss topology encompassing point-wise absolute error, intensity-weighted Sobel gradient consistency, and discrete Laplacian curvature regularization, we hypothesize that the proposed model will significantly outperform classical approaches in noisy, real-world optical regimes.
\end{abstract}

\section{Introduction}
Optical interferometry relies on the superposition of electromagnetic wavefields to extract high-resolution topographical and structural data. However, sensors can only record scalar intensity. Consequently, the true continuous absolute physical phase $\phi(x,y)$ is recorded mathematically wrapped within the principal topological domain.

Let the wrapping operator $\mathcal{W}: \mathbb{R} \to [-\pi, \pi)$ be defined as:
\begin{equation}
  \mathcal{W}(\phi(x,y)) = \phi(x,y) - 2\pi k(x,y)
\end{equation}
where the topological shift field $k(x,y) \in \mathbb{Z}$ represents the unknown integer multiple of $2\pi$ at spatial coordinate $(x,y)$. Recovering the absolute continuous manifold $\phi(x,y)$ from the discrete observation $\mathcal{W}(\phi)$ is an ill-posed inverse problem due to the loss of $k(x,y)$.

Classical phase unwrapping techniques generally fall into two categories: path-following and minimum-norm approaches. Path-following methods attempt to evaluate the discrete gradient $\nabla \mathcal{W}(\phi)$, but struggle catastrophically when high-frequency optical noise triggers phase singularities, leading to unresolvable spatial error propagation. Minimum-norm approaches minimize derivative discrepancies via Poisson-like equations but fundamentally destroy the absolute macroscopic piston ($2\pi k$) structure.

Recent advancements in deep learning have demonstrated the capability of convolutional neural networks (CNNs) to implicitly map the nonlinear relationship between fringe intensity and absolute phase. In this research, we propose a physics-informed deep learning architecture designed to overcome the critical limitations of standard CNNs, specifically the $2\pi k$ translational ambiguity and the lack of physical noise resilience.

\section{Deterministic Synthetic Data Generation and Physical Modeling}

A fundamental vulnerability of deep learning in physical sciences is domain shift caused by insufficient training diversity. To ensure robust generalization to empirical laboratory data while maintaining benchmark integrity, we developed a highly rigorous, deterministic synthetic data generator.

\subsection{The Reproducibility Guarantee}
To generate massive training distributions without compromising scientific benchmarking, the computational pipeline employs a master-seed architecture. A single root cryptographic seed deterministically branches into thousands of completely independent sub-seeds for concurrent multiprocess workers. This mathematically guarantees zero sample overlap or hash collisions across a $100,000+$ image dataset, while ensuring the entire distribution can be perfectly regenerated by a third-party reviewer to validate performance benchmarks.

\subsection{Forward Optical Model}
The synthetic interferograms generated for this study model the interference of two coherent light beams. To standardize the topological geometry of the absolute phase fields, the optical generator strictly utilizes a fixed Helium-Neon laser wavelength of $\lambda = 632.8$ nm.

The primary wavefront $\phi_1(x,y)$ represents structural topography augmented with low-order polynomial aberrations. The reference beam $\phi_2(x,y)$ is mathematically modeled as a continuous divergent spherical wave originating from a longitudinal focal distance of $Z = 0.5$ m:
\begin{equation}
  \phi_2(x,y) = \frac{2\pi}{\lambda} \sqrt{x^2 + y^2 + Z^2}
\end{equation}

Let the base complex amplitude be denoted by $E_0$. The recorded intensity field $I(x,y)$ represents the squared magnitude of the superposition of the two complex wavefields, resulting in the fundamental non-linear interferometry equation:
\begin{equation}
  I(x,y) = \left| E_0 e^{i\phi_1(x,y)} + E_0 e^{i\phi_2(x,y)} \right|^2 = 2 E_0^2 \Big( 1 + \cos(\Delta\phi(x,y)) \Big)
\end{equation}
where $\Delta\phi(x,y) = \phi_1(x,y) - \phi_2(x,y)$. This expansion rigorously demonstrates how the continuous physical phase difference is locked inside a periodic, non-invertible cosine manifold.

\subsection{Supervised Regression Target}
The objective of the model is to map the observed intensity $I(x,y)$ directly to the absolute phase difference between the two beams. To constrain the continuous scalar phase space, the supervised target $\Delta\phi(x,y)$ is shifted to a positive domain:
\begin{equation}
  \Delta\phi(x,y) = (\phi_1 - \phi_2) - \min(\phi_1 - \phi_2)
\end{equation}

\section{Physics-Based Dynamic Noise Curriculum}
Neural networks are highly susceptible to domain shift when applied to empirical laboratory data. To close the gap between synthetic data and physical instrumentation, we introduce an extensive optical degradation pipeline mathematically anchored to the hardware physics of interferometric sensing.

\begin{figure}[H]
  \centering
  \includegraphics[width=\textwidth]{../../results/figs/curriculum_noise_grid.png}
  \caption{Progression of the Hybrid N-Cycle Stochastic Curriculum. The visualization demonstrates the deterministic injection of optical aberrations, laser speckle, and quantum shot noise. By Epoch 35, the network evaluates under catastrophic signal degradation, mimicking severe real-world sensor bounds.}
  \label{fig:curriculum}
\end{figure}

\subsection{Infinite GPU-Native Degradation}
Classical deep learning datasets often compute noisy input pairs offline to mitigate input/output (I/O) bottlenecks. However, this permits the network to eventually memorize specific noise configurations. We circumvent this by porting the entire physical degradation model directly into the GPU batch-preparation thread. This guarantees that at every single epoch, the network encounters a mathematically unique, non-repeating noise topology for an identical underlying absolute phase structure, effectively acting as an infinite dataset.

The noise synthesis adheres strictly to the physical order of operations within an optical sensing array:
\begin{enumerate}
  \item \textbf{Lens Aberration:} Optical defocus modeled as continuous Gaussian blur.
  \item \textbf{Laser Speckle:} Multiplicative coherent noise: $I = I \cdot (1 + \mathcal{N}(0, \sigma_{\text{speckle}}))$.
  \item \textbf{Uneven Illumination:} Low-frequency bicubic beam misalignment.
  \item \textbf{Sensor Electronics:} Photometric gain and black-level offset jitter.
  \item \textbf{Quantum Shot Noise:} Photon arrival statistics mathematically enforced via a positive-domain Poisson counting process. Crucially, this operation is executed prior to any non-physical data standardization (e.g., zero-mean Z-scoring) to preserve strict physical intensity boundaries.
  \item \textbf{Thermal Noise:} Additive Gaussian white noise modeling sensor dark current.
\end{enumerate}

\subsection{Hybrid Two-Cycle Stochastic Curriculum}
Subjecting a neural network to peak optical degradation immediately at initialization typically causes gradient divergence. Conversely, training a network to convergence on a strictly static noise level severely narrows the breadth of the local minima it settles into.

To engineer maximum robustness, we introduce a \textit{Hybrid Two-Cycle Stochastic Curriculum}.
As illustrated in Figure \ref{fig:curriculum}, the network is subjected to an evolving degradation pipeline. Macroscopically, the overall noise level $\mu(\text{epoch})$ does not monotonically increase; instead, it forms a continuous two-cycle cosine wave. During the first cycle, the noise gracefully rises to push the network to its absolute structural limit (e.g., Epoch 35, where the input intensity is dominated by laser speckle and quantum shot noise) before dropping back to near-zero. This forceful variation ejects the model from sharp, brittle local minima. The second cycle then repeats this trajectory, allowing the network to permanently settle into an extremely robust, generalizable global minimum while fine-tuning its final absolute phase mapping on clean, deterministic structures.

Microscopically, the curriculum level is subjected to continuous Gaussian jitter $\gamma_{\text{batch}} \sim \mathcal{N}(\mu, \sigma_{\text{jitter}})$. This ensures that even at the macroscopic peak of the curriculum, the network is randomly forced to correctly reconstruct relatively clean interferograms, preventing catastrophic forgetting of the foundational physical mapping.

\section{Network Architecture and Spatial Anchoring}

The backbone of the proposed method is a highly modified, physics-informed fully convolutional network (UNetRes2) specifically designed to resolve the macroscopic topological ambiguity of phase unwrapping while preserving high-frequency fringe structures.

\subsection{Global Offset Head and Duality Mapping}
The absolute physical phase represents a dual-scale topological space: high-frequency local fringe structures governed by continuous spatial gradients, and a low-frequency macroscopic piston shift ($2\pi k$). To mirror this physical duality, our architecture explicitly bifurcates the regression task. The standard U-Net decoding branch reconstructs the local topographical manifold $\phi_{\text{raw}}(x,y)$. Simultaneously, a parallel convolutional head stems from the deepest bottleneck layer and passes through Global Average Pooling to predict a singular scalar shift $k_{\text{off}}$. The final reconstructed absolute phase is assembled as $\phi_{\text{abs}}(x,y) = \phi_{\text{raw}}(x,y) + k_{\text{off}}$, preventing the network from distorting high-frequency local features to compensate for global scalar errors.

\subsection{Depthwise Separable Res2 Blocks}
Standard $3 \times 3$ convolutions possess insufficient receptive fields to track macroscopic phase trends across high-resolution $256 \times 256$ interferograms. To exponentially expand the spatial context without inducing catastrophic overfitting, the standard U-Net blocks are entirely replaced with Depthwise Separable Res2Net modules. Each block splits the channel dimension into $s=4$ discrete groups and applies sequential, hierarchical depthwise convolutions with cumulative residual connections. This mathematically guarantees a massive effective receptive field capable of anchoring global phase structures while maintaining extreme parameter efficiency.

\subsection{Multi-Scale Deep Supervision}
Phase unwrapping networks frequently suffer from vanishing gradients and structural collapse due to the highly oscillatory nature of the input intensity field. To enforce structural coherence early in the training dynamics, the architecture utilizes Multi-Scale Deep Supervision. Auxiliary regression heads branch off the $1/4\times$ and $1/2\times$ scale decoding layers, returning explicit low-resolution phase predictions. Gradients are injected directly into the intermediate layers of the network during the backward pass, forcing the early decoding blocks to securely establish the macroscopic topological shape before fine-tuning the microscopic details at the terminal $1\times$ resolution.

\subsection{Spatial Anchoring and Phase Hinting}
Standard convolutions are translationally invariant, which fundamentally contradicts the spatially grounded nature of absolute phase unwrapping. Our architecture utilizes an \texttt{AddCoords} module, explicitly concatenating normalized spatial grid coordinates $(x,y) \in [-1, 1]^2$ to the input tensor to break the spatial invariance. Furthermore, the network is provided with a heavily regularized scalar hint $\phi_{\text{hint}}$, corrupted by a random scalar perturbation $\delta$, to force the network to actively align its predicted manifold with physical macroscopic anchors.

\section{Multi-Term Topological Loss}
Due to the ill-posed nature of absolute phase unwrapping, optimizing the network weights $\theta$ exclusively via a point-wise absolute error typically fails to capture the high-frequency structural integrity of the continuous optical surface. Consequently, the network is optimized via a multi-objective loss topology $\mathcal{L}_{\text{total}}$ that rigorously enforces physical continuity, spatial derivatives, and localized confidence weighting:

\begin{equation}
  \mathcal{L}_{\text{total}} = \mathcal{L}_{\text{MAE}} + \mathcal{L}_{\text{grad}} + \mathcal{L}_{\text{curv}}
\end{equation}

\subsection{Global Absolute Accuracy ($\mathcal{L}_{\text{MAE}}$)}
The foundation of the optimization is the global $L_1$ norm (Mean Absolute Error), defined over the spatial dimensions $N = H \times W$. This term strictly anchors the unwrapped manifold to the absolute domain to eliminate the macroscopic $2\pi k$ piston shift:
\begin{equation}
  \mathcal{L}_{\text{MAE}} = w_{\text{mae}} \frac{1}{N} \sum_{x,y} \left| \phi_{\text{pred}}(x,y) - \phi_{\text{gt}}(x,y) \right|
\end{equation}

\subsection{Intensity-Weighted Gradient Consistency ($\mathcal{L}_{\text{grad}}$)}
To enforce high-fidelity structural reconstruction, we compute the spatial depthwise phase derivatives. We approximate the continuous gradients $(\nabla_x, \nabla_y)$ utilizing discrete $3 \times 3$ Sobel convolution kernels, denoted $S_x$ and $S_y$:
\begin{equation}
  \nabla_x \phi \approx S_x \ast \phi, \quad \nabla_y \phi \approx S_y \ast \phi
\end{equation}
Furthermore, we introduce a physical confidence weighting mechanism. We mathematically condition the gradient loss on the square root of the normalized raw intensity $I(x,y)$, effectively masking out derivative errors in dark regions while heavily penalizing structural topological deviations in regions of bright interference fringes:
\begin{align}
  \mathcal{L}_{\text{grad}} = w_{\text{grad}} \frac{1}{N} \sum_{x,y} \sqrt{I(x,y)} \cdot \Big(
    &\left| \nabla_x\phi_{\text{pred}} - \nabla_x\phi_{\text{gt}} \right| + \nonumber \\
  &\left| \nabla_y\phi_{\text{pred}} - \nabla_y\phi_{\text{gt}} \right| \Big)
\end{align}

\subsection{Laplacian Smoothness Regularization ($\mathcal{L}_{\text{curv}}$)}
To enforce optical smoothness and penalize structural artifacting, we apply an $L_1$ penalty to the unrolled predicted surface utilizing a discrete 2D Laplacian operator.
\begin{equation}
  \mathcal{L}_{\text{curv}} = w_{\text{curv}} \frac{1}{N} \sum_{i,j} \left| \nabla^2 \phi_{\text{pred}}(i,j) \right|
\end{equation}

\section{Empirical Results and Evaluation}
Evaluating absolute phase reconstruction is complex due to global scalar shifts (the "piston" effect). Before calculating final diagnostic metrics (RMSE, MAE, SSIM), the validation pipeline performs a Least-Squares Affine Alignment on the raw prediction:
\begin{equation}
  a \cdot \phi_{\text{pred}} + c \approx \phi_{\text{gt}}
\end{equation}
Metrics are computed on the aligned prediction $\hat{\phi} = a\phi_{\text{pred}} + c$. This explicitly isolates the network's topological unwrapping capability from trivial global drift errors.

\subsection{Convergence and Quantitative Metrics}
The network was trained over 100 epochs utilizing a OneCycle learning rate policy. As demonstrated in Figure \ref{fig:training_curve}, the training dynamics exhibit a strong correlation between the learning rate and the validation loss. Despite the aggressive injection of stochastic noise across the curriculum cycles, the Validation Mean Absolute Error (MAE) remains remarkably stable and drops sharply as the learning rate enters its cosine annealing tail. This confirms that the network successfully avoids gradient divergence. Table \ref{tab:metrics} summarizes the final quantitative evaluation. The network achieves a final TopoMAE of 0.5248 rad and a PSNR of 27.44 dB. The sub-radian MAE combined with a high Structural Similarity Index (SSIM = 0.9847) demonstrates the network's capability to reconstruct the continuous phase surface with high physical fidelity.

\begin{table}[H]
  \centering
  \caption{Empirical Validation Metrics at Final Epoch.}
  \label{tab:metrics}
  \begin{tabular}{lrrrrrrr}
    \toprule
    Epoch & TopoMAE & AbsMAE & RMSE & SSIM & PSNR & MaxErr & GradMAE \\
    \midrule
    1 & 4.5673 & 12.3717 & 7.6240 & 0.6486 & 11.28 & 18.447 & 3.3859 \\
    10 & 3.0573 & 5.5060 & 6.2046 & 0.8316 & 13.62 & 13.025 & 1.5016 \\
    20 & 1.1629 & 2.0590 & 4.0077 & 0.9544 & 20.41 & 5.083 & 0.5573 \\
    30 & 1.6164 & 2.6989 & 4.4901 & 0.9251 & 17.33 & 7.027 & 0.7073 \\
    40 & 0.7282 & 1.1633 & 3.1987 & 0.9789 & 25.54 & 3.188 & 0.3510 \\
    50 & 0.7105 & 1.1459 & 3.2076 & 0.9781 & 25.15 & 3.028 & 0.3365 \\
    60 & 0.8147 & 1.5273 & 2.8941 & 0.9682 & 22.49 & 3.690 & 0.3959 \\
    70 & 0.5623 & 0.9985 & 2.8094 & 0.9843 & 28.21 & 2.659 & 0.2780 \\
    80 & 0.4379 & 0.7265 & 2.2191 & 0.9889 & 31.75 & 1.961 & 0.2449 \\
    90 & 0.6569 & 1.1213 & 2.7561 & 0.9762 & 25.20 & 3.070 & 0.3170 \\
    100 & 0.5248 & 0.9818 & 2.5522 & 0.9847 & 27.44 & 2.638 & 0.2599 \\
    \bottomrule
  \end{tabular}
\end{table}

\begin{figure}[H]
  \centering
  \includegraphics[width=\linewidth]{../../results/figs/training_curve.png}
  \caption{Training convergence dynamics. The Validation MAE (orange) rapidly drops in tandem with the learning rate schedule (bottom panel), confirming the absence of gradient divergence despite the stochastic noise curriculum.}
  \label{fig:training_curve}
\end{figure}

\subsection{Spatial Topological Resilience}
To evaluate structural coherence, we visualize the predicted phase manifold against severely corrupted inputs in Figure \ref{fig:qualitative}. The input interferogram (Noisy I) is heavily degraded by 2\% dead-pixel dropouts, dense laser speckle, and optical aberrations, rendering traditional path-following unwrapping algorithms virtually useless. Despite these conditions, the network accurately reconstructs the continuous macroscopic topology (Predicted Abs). The 2D residuals (Error) map shows only minimal structural deviation, primarily localized to regions of extreme noise rather than systematic topographical failures. This demonstrates robustness against severe optical artifacts present in simulated sensor data.
\begin{figure}[H]
  \centering
  \includegraphics[width=\textwidth]{../../results/figs/qualitative_grid.png}
  \caption{Qualitative Spatial Resilience. Despite heavy physical degradation (Noisy I), the predicted absolute phase manifold accurately matches the Ground Truth topology, demonstrating the network's ability to act as a robust physical prior.}
  \label{fig:qualitative}
\end{figure}

\subsection{Sub-Fringe Cross-Sectional Precision}
While 2D topological mapping is critical, quantitative optical metrology requires extreme sub-fringe accuracy along 1D cross-sections. A rigorous 1D cross-sectional analysis over four randomly selected spatial profiles is presented in Figure \ref{fig:phase_profile}. The red dashed line representing the network's prediction flawlessly overlays the black solid line of the ground truth. The 95\% confidence intervals (shaded regions) closely track the ground truth across both horizontal and vertical physical contours, with a maximum residual variance of approximately $\sigma \approx 0.05$ rad. This variance corresponds to less than $1/100$th of a full $2\pi$ fringe, confirming that the network achieves high-confidence sub-fringe precision even when resisting highly oscillatory noise artifacts in the raw intensity.
\begin{figure}[H]
  \centering
  \includegraphics[width=0.48\textwidth]{../../results/figs/phase_profile.png}\hfill
  \includegraphics[width=0.48\textwidth]{../../results/figs/phase_profile_1.png}\\[1em]
  \includegraphics[width=0.48\textwidth]{../../results/figs/phase_profile_2.png}\hfill
  \includegraphics[width=0.48\textwidth]{../../results/figs/phase_profile_3.png}
  \caption{1D Phase Profile Analysis over four random spatial geometries. The network tracks both the horizontal and vertical physical contours with sub-fringe accuracy, resisting highly oscillatory noise artifacts.}
  \label{fig:phase_profile}
\end{figure}

\subsection{Graceful Failure and Outlier Statistics}
Figure \ref{fig:error_hist} characterizes the network's behavior near the mathematical limits of interferometric reconstruction. The error distribution is highly non-Gaussian, featuring a massive central peak and a sparse tail. The low Median MAE ($0.0777$ rad) indicates accurate operation for the majority of the phase space. However, the Mean MAE ($0.5101$ rad) is heavily skewed by 99th-percentile outliers. Crucially, these outliers are tightly clustered around discrete multiples of $2\pi$. This demonstrates a "graceful failure" regime: when the network fails under extreme signal-to-noise ratios, it defaults to classical $2\pi$ topological jumps rather than generating unpredictable, chaotic artifacts that would corrupt the entire reconstructed manifold.
\begin{figure}[H]
  \centering
  \includegraphics[width=\linewidth]{../../results/figs/error_histogram.png}
  \caption{Statistical Distribution of Absolute Phase Errors. The disparity between the Median and Mean MAE highlights the network's resilience, restricting failures to strict $2\pi$ mathematical limits under catastrophic noise.}
  \label{fig:error_hist}
\end{figure}

\section{Future Work}
\subsection{Multi-Seed Statistical Validation}
To account for seed variance, future iterations of this study will deploy an automated multi-seed evaluation. The pipeline will sequentially execute $N=3$ independent training runs utilizing identical configuration limits, each initialized with a different network weight seed. The final TopoMAE and PSNR results will be aggregated to formulate a robust Mean $\pm$ Standard Deviation.

\section{Conclusion}
By integrating deep convolutional learning with established optical physics, the proposed pipeline represents a robust methodology for phase unwrapping. The implementation of spatial coordinate embeddings, dynamic optical noise scheduling, and multi-derivative loss topologies ensures that the network is explicitly trained to handle empirical, hardware-degraded interferograms with high accuracy.

\end{document}
